% Mathematical TeX extracted from the cited web pages, reformatted by the assistant. Not the native full chapter.
\begin{definition}[10.160.4--5]
Let $(R,\mathfrak m)$ be a complete local ring. A subring $\Lambda\subset R$ is called a coefficient ring if the following conditions hold:
\begin{enumerate}
\item $\Lambda$ is a complete local ring with maximal ideal $\Lambda\cap\mathfrak m$;
\item the residue field of $\Lambda$ maps isomorphically to the residue field of $R$; and
\item $\Lambda\cap\mathfrak m=p\Lambda$, where $p$ is the characteristic of the residue field of $R$.
\end{enumerate}
A Cohen ring is a complete discrete valuation ring with uniformizer $p$ a prime number.
\end{definition}

\begin{lemma}[10.159.1]
Let $(R,\mathfrak m,k)$ be a local ring. Let $K/k$ be a field extension. There exists a local ring $(R',\mathfrak m',k')$, a flat local ring map $R\to R'$ such that $\mathfrak m'=\mathfrak mR'$ and such that $k'$ is isomorphic to $K$ as an extension of $k$.
\end{lemma}
\begin{proof}
Suppose that $k'=k(\alpha)$ is a monogenic extension of $k$. Then $k'$ is the residue field of a flat local extension $R\subset R'$ as in the lemma. Namely, if $\alpha$ is transcendental over $k$, then we let $R'$ be the localization of $R[x]$ at the prime $\mathfrak mR[x]$. If $\alpha$ is algebraic with minimal polynomial $T^d+\sum\overline\lambda_iT^{d-i}$, then we let $R'=R[T]/(T^d+\sum\lambda_iT^{d-i})$.

Consider the collection of triples $(k',R\to R',\phi)$, where $k\subset k'\subset K$ is a subfield, $R\to R'$ is a local ring map as in the lemma, and $\phi:R'\to k'$ induces an isomorphism $R'/\mathfrak mR'\cong k'$ of $k$-extensions. These form a ``big'' category $\mathcal C$ with morphisms $(k_1,R_1,\phi_1)\to(k_2,R_2,\phi_2)$ given by ring maps $\psi:R_1\to R_2$ such that
\[
\xymatrix{R_1\ar[d]_\psi\ar[r]_{\phi_1}&k_1\ar[r]&K\ar@{=}[d]\\
R_2\ar[r]^{\phi_2}&k_2\ar[r]&K}
\]
commutes. This implies that $k_1\subset k_2$.

Suppose that $I$ is a directed set, and $((R_i,k_i,\phi_i),\psi_{ii'})$ is a system over $I$, see Categories, Section 4.21. In this case we can consider $R'=\operatorname{colim}_{i\in I}R_i$. This is a local ring with maximal ideal $\mathfrak mR'$, and residue field $k'=\bigcup_{i\in I}k_i$. Moreover, the ring map $R\to R'$ is flat as it is a colimit of flat maps (and tensor products commute with directed colimits). Hence we see that $(R',k',\phi')$ is an ``upper bound'' for the system.

An almost trivial application of Zorn's Lemma would finish the proof if $\mathcal C$ was a set, but it isn't. (Actually, you can make this work by finding a reasonable bound on the cardinals of the local rings occurring.) To get around this problem we choose a well ordering on $K$. For $x\in K$ we let $K(x)$ be the subfield of $K$ generated by all elements of $K$ which are $\leq x$. By transfinite recursion on $x\in K$ we will produce ring maps $R\subset R(x)$ as in the lemma with residue field extension $K(x)/k$. Moreover, by construction we will have that $R(x)$ will contain $R(y)$ for all $y\leq x$. Namely, if $x$ has a predecessor $x'$, then $K(x)=K(x')[x]$ and hence we can let $R(x')\subset R(x)$ be the local ring extension constructed in the first paragraph of the proof. If $x$ does not have a predecessor, then we first set $R'(x)=\operatorname{colim}_{x'<x}R(x')$ as in the third paragraph of the proof. The residue field of $R'(x)$ is $K'(x)=\bigcup_{x'<x}K(x')$. Since $K(x)=K'(x)[x]$ we see that we can use the construction of the first paragraph of the proof to produce $R'(x)\subset R(x)$. This finishes the proof of the lemma.
\end{proof}

\begin{lemma}[10.160.6]
Let $p$ be a prime number. Let $k$ be a field of characteristic $p$. There exists a Cohen ring $\Lambda$ with $\Lambda/p\Lambda\cong k$.
\end{lemma}
\begin{proof}
First note that the $p$-adic integers $\Z_p$ form a Cohen ring for $\mathbf F_p$. Let $k$ be an arbitrary field of characteristic $p$. Let $\Z_p\to R$ be a flat local ring map such that $\mathfrak m_R=pR$ and $R/pR=k$, see Lemma 10.159.1. By Lemma 10.97.5 the completion $\Lambda=R^\wedge$ is Noetherian. Since $\Lambda/p\Lambda=R/pR$ is a field we conclude $\Lambda$ is a complete Noetherian local ring with maximal ideal $(p)$ (equality and completeness by Lemma 10.96.3). For every $n$ the map $\Z/p^n\Z\to\Lambda/p^n\Lambda=R/p^nR$ is flat (equality by Lemma 10.96.3 again). Hence $\Z_p\to\Lambda$ is flat for example by Lemma 10.99.11. Hence $p$ is a nonzerodivisor in $\Lambda$. Hence $\Lambda$ has dimension $\geq1$ (Lemma 10.60.13) and we conclude that $\Lambda$ is regular of dimension $1$, i.e., a discrete valuation ring by Lemma 10.119.7. We conclude $\Lambda$ is a Cohen ring for $k$.
\end{proof}
\begin{lemma}[10.160.7]
Let $p>0$ be a prime. Let $\Lambda$ be a Cohen ring with residue field of characteristic $p$. For every $n\geq1$ the ring map
\[
\Z/p^n\Z\to\Lambda/p^n\Lambda
\]
is formally smooth.
\end{lemma}
\begin{proof}
If $n=1$, this follows from Proposition 10.158.9. For general $n$ we argue by induction on $n$. Namely, if $\Z/p^n\Z\to\Lambda/p^n\Lambda$ is formally smooth, then we can apply Lemma 10.138.12 to the ring map $\Z/p^{n+1}\Z\to\Lambda/p^{n+1}\Lambda$ and the ideal $I=(p^n)\subset\Z/p^{n+1}\Z$.
\end{proof}

\begin{theorem}[10.160.8, Cohen structure theorem]
Let $(R,\mathfrak m)$ be a complete local ring.
\begin{enumerate}
\item $R$ has a coefficient ring (see Definition 10.160.4).
\item If $\mathfrak m$ is a finitely generated ideal, then $R$ is isomorphic to a quotient
\[
\Lambda[[x_1,\ldots,x_n]]/I
\]
where $\Lambda$ is either a field or a Cohen ring.
\end{enumerate}
\end{theorem}
\begin{proof}
Let us prove a coefficient ring exists. First we prove this in case the characteristic of the residue field $\kappa$ is zero. Namely, in this case we will prove by induction on $n>0$ that there exists a section
\[
\varphi_n:\kappa\longrightarrow R/\mathfrak m^n
\]
to the canonical map $R/\mathfrak m^n\to\kappa=R/\mathfrak m$. This is trivial for $n=1$. If $n>1$, let $\varphi_{n-1}$ be given. The field extension $\kappa/\Q$ is formally smooth by Proposition 10.158.9. Hence we can find the dotted arrow in the following diagram
\[
\xymatrix{R/\mathfrak m^{n-1}&R/\mathfrak m^n\ar[l]\\
\kappa\ar[u]^{\varphi_{n-1}}\ar@{..>}[ru]&\Q\ar[l]\ar[u]}
\]
This proves the induction step. Putting these maps together
\[
\varprojlim_n\varphi_n:\kappa\longrightarrow R=\varprojlim_nR/\mathfrak m^n
\]
gives a map whose image is the desired coefficient ring.

Next, we prove the existence of a coefficient ring in the case where the characteristic of the residue field $\kappa$ is $p>0$. Namely, choose a Cohen ring $\Lambda$ with $\kappa=\Lambda/p\Lambda$, see Lemma 10.160.6. In this case we will prove by induction on $n>0$ that there exists a map
\[
\varphi_n:\Lambda/p^n\Lambda\longrightarrow R/\mathfrak m^n
\]
whose composition with the reduction map $R/\mathfrak m^n\to\kappa$ produces the given isomorphism $\Lambda/p\Lambda=\kappa$. This is trivial for $n=1$. If $n>1$, let $\varphi_{n-1}$ be given. The ring map $\Z/p^n\Z\to\Lambda/p^n\Lambda$ is formally smooth by Lemma 10.160.7. Hence we can find the dotted arrow in the following diagram
\[
\xymatrix{R/\mathfrak m^{n-1}&R/\mathfrak m^n\ar[l]\\
\Lambda/p^n\Lambda\ar[u]^{\varphi_{n-1}}\ar@{..>}[ru]&\Z/p^n\Z\ar[l]\ar[u]}
\]
This proves the induction step. Putting these maps together
\[
\varprojlim_n\varphi_n:\Lambda=\varprojlim_n\Lambda/p^n\Lambda\longrightarrow R=\varprojlim_nR/\mathfrak m^n
\]
gives a map whose image is the desired coefficient ring.

The final statement of the theorem follows readily. Namely, if $y_1,\ldots,y_n$ are generators of the ideal $\mathfrak m$, then we can use the map $\Lambda\to R$ just constructed to get a map
\[
\Lambda[[x_1,\ldots,x_n]]\longrightarrow R,\qquad x_i\longmapsto y_i.
\]
Since both sides are $(x_1,\ldots,x_n)$-adically complete this map is surjective by Lemma 10.96.1 as it is surjective modulo $(x_1,\ldots,x_n)$ by construction.
\end{proof}

